\section{Source audit, verification, and further research}
\label{bhr:sec:audit}

\subsection{A precise correction to a published Barnes identity}

The inspection of the Barnes Fourier source revealed a pointwise
algebraic error in equation (7) of the inspected arXiv version of
Mez\H{o}'s paper \cite{bhr:Mezo}. That equation identifies
$\log^2G(x)$ with
\[
 \frac14\left\{
   \log^2\bigl(G(x)G(1-x)\bigr)
   +\log^2\bigl(G(x)/G(1-x)\bigr)\right\}.
\]
On $(0,1)$ both Barnes values are positive. Put
$a=\log G(x)$ and $b=\log G(1-x)$. The displayed expression
is $(a^2+b^2)/2$, not $a^2$. Thus the corrected pointwise identity is
\begin{equation}\label{bhr:eq:Mezo-correction}
 \boxed{\displaystyle
 \frac{\log^2G(x)+\log^2G(1-x)}2
 =\frac14\left\{
   \log^2\bigl(G(x)G(1-x)\bigr)
   +\log^2\bigl(G(x)/G(1-x)\bigr)\right\}.}
\end{equation}
Alternatively, to keep the original left side one must add
\[
 \frac12\log\bigl(G(x)G(1-x)\bigr)
           \log\bigl(G(x)/G(1-x)\bigr)
\]
to the original right side.

This is an exact objection, not a small numerical discrepancy.
From $G(1+x)=\Gamma(x)G(x)$ and $G(1)=1$ one gets
$G(x)=x(1+O(x))$ as $x\downarrow0$, while
$\log G(1-x)=O(x)$. Therefore the original left side has leading
term $\log^2x$, whereas its stated right side has leading term
$\tfrac12\log^2x$.

The distinction has a narrow consequence. After integration over
$(0,1)$, reflection makes
$\int\log^2G(1-x)=\int\log^2G(x)$, and the added product
is odd under $x\mapsto1-x$. The integrated equality used next in
that paper is valid. We therefore propose correcting the pointwise
line or placing its equality after integration; we do not infer
that its subsequent integrated theorem is false. The source version
and location are recorded in \src{integration/CORRECTIONS.md}.

\subsection{What the repository audit does and does not establish}

The source inventory contains 81 manuscript chapter files and all
11 incoming archives present at the pinned commit. The detailed
mathematical review concentrates on the conjecture status records,
the cotangent-resolvent programme, the harmonic power bridges, and
the cyclic Tornheim remainder. The inventory is not a claim that
every theorem in all these documents has received a fresh proof audit.

The following distinctions are necessary for correct integration.
\begin{enumerate}[leftmargin=*,itemsep=5pt]
 \item The classical Barnes/Hurwitz conversion, the log-$G$ Fourier
 coefficients, balanced negapolygamma primitives, beta-integral
 binomial identities, and symmetric Tornheim desingularization keep
 their existing attribution. The new contribution is the proved
 resolvent and coefficient reduction with the stated normalization.
 \item The earlier report already contains the quartic directional
 Tornheim reconstruction and the formal cyclic dimension count.
 Section~\ref{bhr:fcy:section} supplies the requested explicit convergent
 integral, and then the fifth-order formula using that analytic
 coordinate; it does not relabel the earlier formal result as new.
 \item The quartic digamma moment in
 Section~\ref{bhr:quartic-harmonic:sec} is a different question from
 the quartic directional Tornheim layer. Its depth-three coefficient
 $\beta_1$ is strict, as specified in
 \eqref{bhr:quartic-harmonic:beta}. An inclusive harmonic coefficient
 cannot be substituted without its diagonal correction.
 \item At a resonance, a Laurent constant in the spectral or kernel
 exponent is not automatically the finite part in the original
 integration coordinate. Equations~\eqref{bhr:eq:weighted-contact}
 and \eqref{bhr:cxp:eq:lambdaCT} explicitly perform the conversion.
 \item The current $S_6$ and revised $S_8$ coefficient vectors remain
 conjectural. The pre-existing rejection of the old frozen $S_8$
 vector must retain its separate label and status. No present
 decimal evaluation changes any of these statuses.
\end{enumerate}

Within the inspected repository material we found no additional
mathematical error that is established by the present analysis.
The pointwise Barnes correction above is to the cited external
source. The attached claim ledger distinguishes proved identities,
classical ingredients, numerical diagnostics, and open questions.
The package consists of research documents and reproducible
calculations, not machine-checked ProveIt proof objects.

\subsection{Independent checks and reproducibility}

Four independent lines of calculation accompany the proofs.
The Barnes checks evaluate the depth-two derivative from its own
double series and evaluate the Fourier sum $\sum q^nT_n$ separately.
The same code compares the complete transform with direct
$\log G$ quadrature in both half-planes, and checks the removable
$q=0$ real component. The polynomial Hurwitz checks compare the
subtracted double series and its spectral derivative with direct
Hurwitz quadrature, including a weighted Stieltjes coefficient.

The complex-power checks compare its Mellin generator with the
original integral at nonintegral complex parameters, and independently
test its Gamma closure and a noncommuting finite-part example.
The quartic checks use three routes: local digamma subtraction,
the convergent harmonic series with a zeta-derivative tail, and
the independent Mellin coefficients $\eta,\beta_1$.
The cyclic fifth checks compare the new ordinary-integral coordinate
with an independently evaluated continuation of three asymmetric
cycles. The pre-existing evaluator used in that comparison is
preserved with its provenance.

Finite exact algebra checks cover the multiple-Gamma coefficient
polynomials and normalization, quartic principal parts and their
finite summation identity, the elementary Tornheim correction,
and the symbolic fifth-order polynomial. These calculations are
useful for detecting signs, factorials, and omitted constants.
The analytic assertions---normal convergence, uniqueness of an
entire remainder, branches, continuation, and the permitted
coefficient interchanges---are proved in the text.

All numerical records are explicitly non-interval diagnostics.
Their working precision and truncations are recorded. Several
independent values agreeing to many digits cannot certify an
integer relation or a finite closed form. Scripts, dependency
versions, replay commands, numerical records, and a file-hash
manifest are included so that integration can preserve both the
theorem and its actual verification status.

\subsection{Further research questions}

The next problems are formulated as exact identities or reductions
in specified function classes. The first four build most directly
on the new formulas.

\begin{enumerate}[label=\textbf{F\arabic*.},leftmargin=*,itemsep=9pt]
 \item \textbf{Products of individual Barnes logarithms.}
 Determine the minimal explicit nested-series class for
 \[
  \int_0^1 W(x)\log\Gamma_r(x)\log\Gamma_t(x)R_z(x)\,dx.
 \]
 The polynomial Hurwitz decomposition reduces this to a bilinear
 weighted transform. The real question is whether its lattice
 splitting closes at depth three, with all diagonal corrections
 explicit, and which reflection combinations reduce the depth.

 \item \textbf{A second representation of $\mathcal D(q)$.}
 Find a Gamma, Barnes, or iterated-integral expression for
 \eqref{bhr:eq:Dq} that is independent of its defining depth-two sum.
 A useful target is an exact differential or functional equation
 under $q\mapsto q/(q-1)$, with a specified branch domain and
 initial value. Establish whether it reduces to ordinary
 polylogarithms and their order derivatives; do not infer such a
 reduction from a successful numerical fit at isolated points.

 \item \textbf{All higher digamma powers.}
 Starting from the principal parts of $\psi(z)^5$, derive the
 exact finite difference of its residue polynomial. Determine
 which height-one depth-four Laurent coefficients and which
 additional harmonic numerators survive summation by parts.
 A satisfactory result must state a finite family of retained
 coefficients and fix the entire remainder and every finite-part
 constant, as was done for the fourth power here.

 \item \textbf{Shifted quartic harmonic bridges.}
 For $a>0$, derive a common-shift analogue of
 \eqref{bhr:quartic-harmonic:main-eq}, using prefixes
 $\sum_{k=0}^{n-1}(k+a)^{-r}$ and the corresponding outer
 denominator. Track simultaneously the origin of the integration
 coordinate, the lower summation endpoint, and generalized
 Stieltjes constants $\gamma_j(a)$. Merely shifting the outer
 denominator leaves a different problem.

 \item \textbf{Nonpositive kernel resonances.}
 Extend the pointwise finite-part comparison of
 Theorem~\ref{bhr:cxp:thm:contacts} to $\Re\lambda\le0$,
 keeping contributions from both endpoints. There is already a
 concrete additional phase at $\lambda=-1$. The continued mean
 of the kernel is $A^\lambda$, so its value there is
 $-z-\epsilon\pi i$, whereas
 \[
  \FP_x\int_0^1 R_z(x)^{-1}\,dx
   =\FP_x\int_0^1(K(x)-z)\,dx=-z.
 \]
 The last equality follows by integrating $K=(\log\sin\pi x)'$
 on the symmetric cutoff interval. Thus the positive-resonance
 rule cannot simply be copied to negative powers.

 \item \textbf{Rational kernel exponents.}
 For a fixed nonintegral rational $\lambda$, determine whether
 $\mathcal L_\lambda(u,q)$ admits a finite representation in a
 specified hypergeometric or multiple-polylogarithm class at
 integral $u$. Its Gamma specialization at $q=0$ and its exact
 exponent jets are constraints on any proposed reduction. The
 unrestricted complex-$\lambda$ Dirichlet series is presently
 retained as an explicit coordinate.

 \item \textbf{The arithmetic of $\beta_1$.}
 Decide whether the normalized integral in
 \eqref{bhr:quartic-harmonic:beta-mellin} reduces to a specified
 algebra generated by ordinary zeta derivatives, Stieltjes
 constants, and Gamma or Barnes values at rational arguments.
 The proved identity for $Q_4+4\gamma Q_3$ gives an exact
 equivalence with a digamma finite-part evaluation, not an
 arithmetic independence statement. A new reduction must add
 information beyond that equivalence.

 \item \textbf{Higher symmetric Tornheim coefficients.}
 Use \eqref{bhr:fcy:Jintegral} to extract the three homogeneous cubic
 coefficients of $J_3$, including their elementary corrections.
 They control the cyclic sixth-derivative layer. Determine which
 reflection or differential identities reduce the formal
 three-coordinate count. For the present quadratic layer, seek
 a second exact representation of $\chi$ that could lead to a
 finite evaluation of $\chi$ or $\Omega_5$.

 \item \textbf{Commensurable Barnes periods.}
 Extend the individual resolvent theorem from unit periods to
 positive rationally commensurable periods. Multiplicity
 sequences should split into finitely many polynomial
 residue-class pieces, suggesting shifted Hurwitz and root-of-unity
 polylogarithm transforms. Prove that reduction with the exact
 normalization; arbitrary incommensurable periods are a separate
 problem and need not remain at finite depth two.

 \item \textbf{Rational shifts and cyclotomic colors.}
 Combine the polynomial weight with a rational Hurwitz shift,
 or with a finite Fourier character. Determine the exact
 colored analogue of $\mathcal O_{\alpha,\beta}(q)$ and its
 residue cancellation. An effective theorem should specify
 cumulative color products and convergence, since individual
 arguments of modulus greater than one need not obstruct the
 coupled nested sum.

 \item \textbf{Exact certificates for the current $S_6$ and $S_8$.}
 Keep each current candidate vector fixed. Try to express its
 difference from the claimed integral in a finite basis of
 iterated integrals or a rigorously defined cyclotomic quotient,
 and reduce that expression by proved shuffle, stuffle,
 distribution, and path relations. The deliverable should be
 an exact certificate, or an exact obstruction to the specified
 reduction, rather than additional matching digits. The
 already-rejected old $S_8$ vector is not the current target.

 \item \textbf{Formal finite-part and convergence infrastructure.}
 Formalize a reusable theorem that subtracts a parameter-dependent
 endpoint jet, integrates its monomials, and compares coefficient
 extraction with the original-coordinate finite part. Pair it
 with normal-convergence lemmas for a geometric outer index and
 a subtracted power-law inner index. These two modules would
 support most identities here without treating each moving pole
 as a separate ad hoc calculation.
\end{enumerate}

\subsection{Integration guide}

The suggested integration order is the polynomial transform,
the individual Barnes theorem, the complex-power generator,
the quartic harmonic bridge, and finally the transverse Tornheim
generator. The first two share their new analytic input;
the remaining branches have their own local subtraction proofs.
All manuscript labels in the deliverable use a common prefix.
The standalone source and the modular source compile the same
article. No external PDFs are bundled: the bibliography links to
the primary sources, while the provenance records identify the
exact repository snapshot and the reused numerical evaluator.
